\documentclass[12pt]{article}

% Page geometry
\usepackage[margin=1in]{geometry}

% Packages for math
\usepackage{amsmath, amsfonts, amssymb}

\begin{document}

\title{Q5: ELCT 562}
\author{Jacob Vaught}
\date{}
\maketitle

\section*{Answer}

\begin{enumerate}
    \item \textbf{Guard Band for the Reconstruction Filter:}

    The guard band in an IQ modulator setup is essential for preventing aliasing and ensuring signal integrity. For DACs with a sampling rate of 20 Msps, the Nyquist frequency is 10 MHz. This frequency represents the upper limit for the signal spectrum in the digital-to-analog conversion process without incurring aliasing. The desired signal bandwidth is 5 MHz (from \( f_c - 2.5 \) MHz to \( f_c + 2.5 \) MHz). The guard band is calculated as the difference between the Nyquist frequency and the highest frequency in the signal bandwidth. Assuming \( f_c \) is chosen to maximize the available bandwidth without aliasing, the highest signal frequency is slightly less than 10 MHz. Therefore, the guard band is the small frequency range between this highest signal frequency and 10 MHz.

    \item \textbf{Example Expression for the Discrete Set of Samples:}

    A sinusoidal waveform can be used as a basic example to demonstrate a signal within the 5 MHz bandwidth. The discrete-time expression for such a signal sampled at 20 Msps is:
    \[ x[n] = A \cdot \sin(2 \pi f n T_s) \]
    Here, \( A \) represents the amplitude of the sinusoid. The frequency \( f \) should be chosen within the 5 MHz bandwidth. The sample index \( n \) iterates over each discrete time point. \( T_s \), the sampling period, is the reciprocal of the sampling frequency, which is \( \frac{1}{20 \text{ MHz}} \). This expression captures the discrete nature of the signal in the time domain, suitable for digital processing and subsequent DAC conversion.

    \item \textbf{Maximum Signal Bandwidth in the Passband:}

    The maximum signal bandwidth that can be synthesized in the passband using this IQ modulator setup is constrained by the Nyquist frequency, which in this case is 10 MHz. This means that the highest frequency that can be accurately represented is 10 MHz. However, for a signal centered around a carrier frequency \( f_c \), with a passband of \( f_c \pm 2.5 \) MHz, the actual bandwidth of interest is 5 MHz. This bandwidth fits within the 10 MHz limit set by the Nyquist criterion, allowing for accurate representation and synthesis of signals in this range. It is important to note that exceeding this bandwidth would result in aliasing, where high-frequency components of the signal are misrepresented as lower-frequency components.
\end{enumerate}

\end{document}
